problem:  
Let \(M^k \subset \mathbb{H}^{n+1}\) be a complete, properly embedded minimal submanifold with smooth, closed asymptotic boundary \(\Lambda = \partial_\infty M \subset S^n_\infty\).  
In the conformal compactification \(g_{\mathbb{H}^{n+1}} = \rho^{-2}(d\rho^2 + h_{S^n})\), write \(M\) as a normal graph over its asymptotic cone \(C_\Lambda\) so that in local coordinates near \(\Lambda\),
\[
u(\theta,\rho) = \rho\, u_1(\theta) + \rho^2 u_2(\theta) + \cdots .
\]
The Graham–Witten expansion expresses the coefficients \(u_j(\theta)\) as tensors built from the induced metric \(h_0\) on \(\Lambda\), its intrinsic curvature, and the second fundamental form \(A_\Lambda\) of the embedding \(\Lambda\subset S^n_\infty\).

**Refined Concrete Problem (Conjecture):**  
The first few coefficients \(u_j\) (\(j\!\le\!n\)) are individually known for small \(n\), and explicit formulas show where \(A_\Lambda\) appears. However, for general codimension and arbitrary \(k \geq 2\), a complete identification of the *first occurrence* of the trace‑free second fundamental form of \(\Lambda\) in the expansion is not known.

- Determine the smallest integer \(j^* = j^*(k,n)\) such that, in the formal expansion of any minimal submanifold \(M^k \subset \mathbb{H}^{n+1}\) with boundary \(\Lambda\), the coefficient \(u_{j^*}\) necessarily contains an **irreducible tensor built from \(A^\circ_\Lambda\)** (the trace-free part of \(A_\Lambda\)) that cannot be recast purely in terms of intrinsic curvature invariants of \(h_0\).  
- In particular, verify whether total umbilicity of \(\Lambda\) implies that all coefficients \(u_j\) for \(j \le j^*\) become intrinsic (functions only of \(h_0\) and its curvature tensors), and identify the order \(j^*\) explicitly.

Why is it a "good" problem:  
This version is sharply defined, grounded in existing asymptotic theory yet pushes beyond known results. It requests a *precise structural threshold*: the first order in the Graham–Witten expansion where extrinsic energy of the boundary embedding irreducibly enters.  

Its significance and depth include:  

- **Bridge between geometry and analysis:** It ties the analytic asymptotic expansion of nonlinear PDEs for minimal submanifolds to algebraic–geometric data of conformal embeddings, linking conformal geometry and minimal surface analysis in a new quantitative way.  
- **Filling a conceptual gap:** For \(k=2\) and low codimension, some coefficients have been computed; for general \(k,n\), the pattern of where purely intrinsic dependence breaks down remains unknown. Identifying \(j^*\) would clarify the divide between “holographic” intrinsic invariants and genuinely extrinsic terms.  
- **Mathematical realism:** The conjecture acknowledges known results (appearance of \(A_\Lambda\) at certain orders) but pushes for a universal geometric characterization of the exact transition point, which is presently not determined by any existing theorem.  
- **Simplicity with depth:** Its statement is short and comprehensible—“At which order does extrinsic curvature first appear?”—yet its resolution would require deep computation in nonlinear geometric analysis, representation theory of curvature tensors, and conformal geometry.

This refined problem is therefore concrete, plausible, mathematically precise, and conceptually rich—qualifying as a truly “good” research‑level mathematical problem.